\par\medskip
\noindent\textbf{Definición 5.9.}
Sea $S$ un esquema formal noetheriano y
$f:\mathfrak X\longrightarrow S$ un esquema formal $S$-ádico
esencialmente de tipo finito. Supongamos que
$J_{\mathfrak X}^p(\mathfrak X|S)$ es propio sobre $S$, de modo que
\[
f_*\mathcal O_{J_{\mathfrak X}^p(\mathfrak X|S)}
\]
es un $\mathcal O_S$-módulo coherente. Definamos
\[
\underline I_S^p(\mathfrak X|S)
=\underline I_S^0\!\left(
 f_*\mathcal O_{J_{\mathfrak X}^p(\mathfrak X|S)}
\right),
\]
y sea $J_S^p(\mathfrak X|S)$ el subesquema formal cerrado de $S$
definido por $\underline I_S^p(\mathfrak X|S)$.

Si existe un entero fijo $n$ que satisface la condición 5.4(1) o
5.4(2) para todo $X_\nu\longrightarrow S_\nu$, pongamos
\[
\begin{aligned}
\underline I_{\mathfrak X}(\mathfrak X|S)
 &=\underline I_{\mathfrak X}^n(\mathfrak X|S),
&\qquad
J_{\mathfrak X}(\mathfrak X|S)
 &=J_{\mathfrak X}^n(\mathfrak X|S),\\
\underline I_S(\mathfrak X|S)
 &=\underline I_S^n(\mathfrak X|S),
&
J_S(\mathfrak X|S)
 &=J_S^n(\mathfrak X|S).
\end{aligned}
\]

\medskip
\noindent\textbf{Corolario 5.10.}
Sea $S$ un esquema formal noetheriano y
$f:\mathfrak X\longrightarrow S$ un esquema formal $S$-ádico
esencialmente de tipo finito tal que
$J_{\mathfrak X}^p(\mathfrak X|S)$ es propio sobre $S$. Entonces:

\smallskip
\noindent\textbf{(1)}
\[
\bigl|J_S^p(\mathfrak X|S)\bigr|
=f\bigl(\,|J_{\mathfrak X}^p(\mathfrak X|S)|\,\bigr).
\]

%%%%%%%%%%  scan idx 134  |  folio 123  %%%%%%%%%%

\smallskip
\noindent\textbf{(2)} La formación de $J_S^p(\mathfrak X|S)$
conmuta con el cambio de base: si
$\mathfrak X'=\mathfrak X\times_SS'$, entonces
\[
J_{S'}^p(\mathfrak X'|S')
=J_S^p(\mathfrak X|S)\times_SS'.
\]

\smallskip
\noindent\textbf{(3)} Si
\[
J_{\mathfrak X}^p(\mathfrak X|S)
=\coprod_{1\leq i\leq m}\mathfrak Y^{(i)}
\]
es una unión disjunta, entonces
\[
\underline I_S^p(\mathfrak X|S)
=\prod_{1\leq i\leq m}
 \underline I_S^0\bigl(f_*\mathcal O_{\mathfrak Y(i)}\bigr).
\]
En particular, si
\[
\bigl|J_{X_0}^p(X_0|S_0)\bigr|=\{x_1,\ldots,x_m\}
\]
es discreto, entonces
\[
\underline I_S^p(\mathfrak X|S)
=\prod_{1\leq i\leq m}
 \underline I_S^p(\mathfrak X_{(x_i)}|S)
=\prod_{1\leq i\leq m}
 \underline I_S^p(\mathfrak X_{(\widehat{x_i})}|S).
\]

\smallskip
\noindent\textbf{(4)} Sea $S$ un esquema noetheriano ordinario,
$X$ un $S$-esquema esencialmente de tipo finito y
$S_0\hookrightarrow S$ un subesquema cerrado. Sea $\widehat S$ la
completación formal de $S$ a lo largo de $S_0$. Si
\[
\widehat X=X\times_S\widehat S
\]
es la completación formal de $X$ a lo largo de
$X_0=X\times_SS_0$, entonces
\[
J_{\widehat S}^p(\widehat X|\widehat S)
=J_S^p(X|S)\times_S\widehat S.
\]
Así, es la completación formal de $J_S^p(X|S)$ a lo largo de
$J_{S_0}^p(X_0|S_0)$.

\smallskip
\noindent\textit{Demostración.}
Por 5.1(c),
\[
\operatorname{Supp}
 f_*\mathcal O_{J_{\mathfrak X}^p(\mathfrak X|S)}
=\bigl|J_S^p(\mathfrak X|S)\bigr|,
\]
lo cual demuestra (1). Las afirmaciones (2) y (4) se deducen de
5.1(a), mientras que (3) se deduce de 5.1(d) y, para un punto aislado
$x\in J_{\mathfrak X}^p(\mathfrak X|S)$, del isomorfismo
\[
J_{\mathfrak X_{(\widehat{x})}}^p
 (\mathfrak X_{(\widehat{x})}|S)
\xrightarrow{\ \sim\ }
J_{\mathfrak X_{(x)}}^p(\mathfrak X_{(x)}|S).
\]
\hfill$\square$

\medskip
\noindent\textbf{Observación 5.11.}
Nuestra definición del ideal jacobiano, adoptada de la definición
clásica sobre un cuerpo base, conmuta con el cambio de base y, bajo la
hipótesis leve de 5.4, proporciona un criterio de lisura. Sin embargo,
en un marco general probablemente se pueda definir de muchas maneras
un subesquema canónico cuyo soporte sea el lugar no liso. Esbozamos una
de esas definiciones.

Para un $A$-álgebra $B$, pongamos
\[
I_{B|A}=\bigcap_E\operatorname{Ann}D^1(B|A,E),
\]
donde $E$ recorre todos los $B$-módulos de tipo finito. Por 3.13(c),
el ideal $I_{B|A}$ conmuta con la localización. Así, un morfismo
$f:X\longrightarrow Y$ determina un haz de ideales
$\underline I_{X|Y}$. Si $f$ es esencialmente de tipo finito, entonces
por 3.13(a), $f$ es liso en $x$ si y solo si
\[
x\notin\operatorname{Supp}
 \bigl(\mathcal O_X/\underline I_{X|Y}\bigr).
\]

%%%%%%%%%%  scan idx 135  |  folio 124  %%%%%%%%%%

Esta construcción se extiende naturalmente a los esquemas formales
ádicos esencialmente de tipo finito y, por tanto, es igualmente
adecuada para la teoría de deformaciones.
